\section{Summary}\label{sec:summary}

In this paper, we have  studied  the quasi-PDF defined with $\gamma^t$ which is free from mixing at $\mathcal{O}(a^0)$. We  have used $M_\pi=356$~MeV Lattice data to demonstrate the matching procedure and show that the excited state contamination is well under control. The one-loop matching coefficient is calculated and we have  discussed the sources of systematic errors as well as the choice of the projection in detail.

We have found  that the systematic uncertainties from the FT truncation method and also the $p_z^R$ dependence are sizable. But  those uncertainties  from $\mu_R$, inversion of matching and choice of projection are relatively minor with $p_z^R$=0. At the same time, the significant change from quasi-PDF to matched PDF suggests that  higher-loop corrections are needed as exhibited in Fig.~\ref{fig:matchingPDF}.


Controlling   systematic uncertainty from the excited state is very challenging since the relative uncertainty grows very fast when either source-sink separation $t_\text{seq}$ or nucleon momentum $P_z$ become large. The two-state fit with smaller separation provides a possibility to obtain a precise result in small $t_\text{seq}<1$fm  region, while for an accurate measurement at large separation using very high statistics, estimating the systematic uncertainty of such a fit is still needed.

Besides the uncertainties that we have studied, in  the future we plan to investigate other systematics such as Lattice discretization and finite volume effects~\cite{Lin:2019ocg} as well as higher twist contributions that affect the small-$x$ result. The latter can be improved with larger nucleon momentum and estimated by extrapolating to infinite nucleon momentum.

Our final result for lightcone PDF agrees with the  global analysis in the large-$x$ region, which gives an encouraging signal  that   LaMET may allow  us to  precisely access  parton physics in the future.


